% 编译配方: xelatex -> bibtex -> xelatex -> xelatex
% minted 代码块需要 -shell-escape。
\documentclass[lang=cn,a4paper,newtx,code=minted,bibtex]{elegantpaper}

\title{Monge 性与在线决策单调性 DP 优化：\\从 SMAWK 到 LARSCH 及其在算法竞赛中的应用}
\author{杨晨 \and 25210980120}
\institute{算法课程期末综述论文}

% \version{修订稿}
\date{\zhdate{2026/06/20}}

\usepackage{array}
\usepackage{tikz}
\usepackage{float}
\usetikzlibrary{arrows.meta,positioning}

\DeclareMathOperator*{\argmin}{arg\,min}
\DeclareMathOperator{\opt}{opt}
\newcommand{\BigO}{\mathcal{O}}
\newcommand{\dd}{\mathrm{d}}
\newcommand{\eps}{\varepsilon}
\newcommand{\monge}{\mathsf{Monge}}

\addbibresource[location=local]{reference.bib}
\AtBeginBibliography{\footnotesize\raggedright}
\setlength{\bibitemsep}{-0.35\baselineskip}

\begin{document}

\maketitle

\begin{abstract}
在近期 AtCoder、Codeforces 与洛谷的若干题解和模板文章中，一个反复出现的现象是：不少实现范式并非全新算法，而是经典矩阵搜索与动态规划加速理论在竞赛约束下的再组织。Monge 性、totally monotone 矩阵、SMAWK 算法与 LARSCH 算法构成了这样一条技术脉络：它们把形如 $dp_i=\min_{j<i}\{dp_j+w(j,i)\}$ 的二次动态规划降到近线性甚至线性复杂度。内容围绕 Monge 矩阵、全单调矩阵与决策单调性的定义，Knuth-Yao 优化、SMAWK 与在线动态规划的代表性结果，以及一种更容易写进竞赛代码的简化 LARSCH 分治框架展开；简化框架包含明确的算法契约、正确性证明和复杂度分析。应用部分结合 AtCoder、Codeforces 与洛谷近年的题解和教程，讨论 WQS 二分、滑动窗口代价、min-plus/max-plus 卷积与背包优化中的使用方式，并以 ABC355G Baseball 为例展示从题意建模、Monge 性证明到实现选择的完整流程。
\keywords{算法竞赛，动态规划优化，Monge 矩阵，SMAWK，LARSCH，决策单调性}
\end{abstract}


\section{引言}

算法竞赛选手常说的近期技巧大致分为两类。一类是某个题目独有的构造、计数或数据结构技巧；另一类则是经典算法理论在竞赛约束下的再包装。洛谷、Codeforces 与 AtCoder 的若干近期题解中出现的 Monge DP、SMAWK、LARSCH、Aliens trick 与 min-plus convolution，大多属于后一类：理论根基较早，但可写性、适用边界与工程组合方式正在变化。

核心问题是一维在线动态规划
\begin{equation}
  dp_i=\min_{0\le j<i}\{dp_j+w(j,i)\},\qquad i=1,2,\ldots,n .
  \label{eq:basic-dp}
\end{equation}
朴素计算需要枚举所有二元组 $(j,i)$，复杂度为 $\BigO(n^2)$。如果代价函数 $w$ 满足 Monge 性，则最优决策点通常具有单调结构，可以使用分治优化、SMAWK 或 LARSCH 等方法降低复杂度。这条技术线索同时具备清楚的数学定义、完整的历史脉络和鲜明的竞赛实践背景。Monge 性可以用一个四点不等式刻画，SMAWK 研究的是全单调矩阵行最值搜索，二者之间存在严格关系；Knuth 对最优二叉搜索树的二次优化、Yao 的四边形不等式框架、Aggarwal 等人的 SMAWK、Larmore 与 Schieber 的在线动态规划，构成了一条自然的算法发展线；AtCoder ABC355 G 的官方题解把分治、SMAWK、LARSCH 与凸包优化并列比较\cite{atcoder2024abc355g}；Codeforces 上也有 SMAWK 作为分治优化替代方案的教程与讨论\cite{codeforces2023smawk}；洛谷 2025 年的算法文章进一步强调简化 LARSCH 在滑动窗口代价函数和莫队代价函数下的实用性\cite{luogu2025larsch}。

结构安排按“定义—历史—算法—应用—现状”展开。第 \ref{sec:def} 节定义 Monge 性、全单调矩阵与在线 DP。第 \ref{sec:history} 节回顾历史代表结果。第 \ref{sec:algorithm} 节给出简化 LARSCH 框架，证明正确性并分析复杂度。第 \ref{sec:applications} 节总结算法竞赛中的应用。第 \ref{sec:sota} 节讨论目前的实践状态、局限和解释框架。

\section{问题定义与基本结构}
\label{sec:def}

\subsection{Monge 矩阵与全单调矩阵}

\begin{definition}[Monge 矩阵]
对一个矩阵 $A$，如果对任意 $i_1<i_2$ 与 $j_1<j_2$ 都有
\begin{equation}
  A_{i_1,j_1}+A_{i_2,j_2}\le A_{i_1,j_2}+A_{i_2,j_1},
  \label{eq:monge}
\end{equation}
则称 $A$ 是用于求最小值的 Monge 矩阵。若不等号方向相反，则对应求最大值版本。
\end{definition}

\begin{definition}[全单调矩阵]
如果矩阵 $A$ 的任意子矩阵中，逐行最小值所在列的下标都是非降的，则称 $A$ 是 totally monotone 矩阵。若只要求原矩阵本身的逐行最小值列非降，则称为单调矩阵。除非特别说明，行最小值均按统一的最左平局规则选取。
\end{definition}

Monge 性比全单调性更容易从具体代价函数中证明。经典结论是：用于求最小值的 Monge 矩阵一定是 totally monotone。SMAWK 的重要性在于，它不需要显式存储矩阵；对 $m\times n$ 的全单调矩阵，只要可以在需要时计算 $A_{i,j}$，就能用 $\BigO(m+n)$ 次矩阵查询求出所有行最值\cite{aggarwal1987geometric}。

SMAWK 可以理解为两个阶段的递归矩阵搜索：先利用全单调性删去不可能成为行最小值的列，使列数降到与行数同阶；再递归求奇数行或偶数行的最小值，并在相邻已知行之间插值搜索剩余行。这个 reduce/interpolate 思路是后续 Monge DP 优化的离线基准。

\begin{theorem}[Monge 蕴含决策单调]
设 $A$ 满足式 \eqref{eq:monge}，并令 $p_i$ 为第 $i$ 行最左的最小值所在列，则 $p_i$ 关于 $i$ 非降。
\end{theorem}

\begin{proof}
反设存在 $i_1<i_2$ 但 $p_{i_1}>p_{i_2}$。令 $j_1=p_{i_2}$，$j_2=p_{i_1}$，于是 $j_1<j_2$。由 Monge 不等式，
\[
  A_{i_1,j_1}+A_{i_2,j_2}\le A_{i_1,j_2}+A_{i_2,j_1}.
\]
另一方面，$j_2$ 是第 $i_1$ 行的最左最小值，$j_1$ 是第 $i_2$ 行的最左最小值，因此
\[
  A_{i_1,j_2}\le A_{i_1,j_1},\qquad A_{i_2,j_1}\le A_{i_2,j_2}.
\]
两组不等式相加只能全部取等。若全部相等，则 $j_1$ 也是第 $i_1$ 行的最小值，且 $j_1<j_2=p_{i_1}$，这与 $p_{i_1}$ 是最左最小值矛盾。因此 $p_{i_1}\le p_{i_2}$。
\end{proof}

任意子矩阵仍满足同一个 Monge 不等式，所以同一证明也给出 totally monotone 性；区别只在于把行、列限制到子矩阵后重新应用最左平局规则。

\subsection{从矩阵搜索到动态规划}

对式 \eqref{eq:basic-dp}，定义合法单元上的隐式矩阵
\begin{equation}
  A_{i,j}=dp_j+w(j,i),\qquad 0\le j<i\le n .
  \label{eq:implicit-matrix}
\end{equation}
求 $dp_i$ 就是在第 $i$ 行的合法列 $j<i$ 中找最小值。非法位置 $j\ge i$ 不参与转移；如果把它们形式上填为 $+\infty$，得到的是阶梯形矩阵，而不是普通的满矩形矩阵。后续算法实际需要的是候选前缀上的阶梯全单调性：固定候选前缀后，按统一平局规则选出的行最小值列随行号非降。

\begin{lemma}[列平移不破坏 Monge 性]
若代价函数满足
\[
  w(a,c)+w(b,d)\le w(a,d)+w(b,c),\qquad a<b<c<d,
\]
则矩阵 $A_{i,j}=dp_j+w(j,i)$ 在所有四个单元均合法的二行二列上仍满足 Monge 不等式。
\end{lemma}

\begin{proof}
取 $j_1<j_2<i_1<i_2$。由 $w$ 的 Monge 性，
\[
  w(j_1,i_1)+w(j_2,i_2)\le w(j_1,i_2)+w(j_2,i_1).
\]
两边同时加上 $dp_{j_1}+dp_{j_2}$，得到
\[
  A_{i_1,j_1}+A_{i_2,j_2}\le A_{i_1,j_2}+A_{i_2,j_1}.
\]
列项 $dp_j$ 在二行二列比较中完全抵消，因此不会改变 Monge 性。合法域的阶梯形限制只减少可比较的单元；实际算法从不查询非法位置。
\end{proof}

\begin{lemma}[候选前缀上的阶梯全单调性]
\label{lem:staircase-prefix}
在式 \eqref{eq:implicit-matrix} 的合法阶梯区域上，若所有合法二行二列都满足 Monge 不等式，并且每行平局时取最左最小列，则对任意候选前缀 $t$，
\[
  q_i^{(t)}=\argmin_{0\le j\le t,\,j<i} A_{i,j}
\]
随行号 $i$ 非降。
\end{lemma}

\begin{proof}
取 $i_1<i_2$，并设 $x=q_{i_1}^{(t)}$、$y=q_{i_2}^{(t)}$。若 $x>y$，则 $y<x<i_1<i_2$，所以 $x,y$ 对两行都合法。由 Monge 不等式，
\[
  A_{i_1,y}+A_{i_2,x}\le A_{i_1,x}+A_{i_2,y}.
\]
又因为 $x$ 是第 $i_1$ 行在该前缀中的最左最小列，$y$ 是第 $i_2$ 行在该前缀中的最左最小列，有
\[
  A_{i_1,x}\le A_{i_1,y},\qquad A_{i_2,y}\le A_{i_2,x}.
\]
两组不等式只能全部取等。于是 $y$ 也是第 $i_1$ 行的最小列，并且 $y<x$，与 $x$ 的最左性矛盾。因此 $x\le y$。若 $y\ge i_1$，则本来就有 $x<i_1\le y$，也不会出现逆序。结论成立。
\end{proof}

这两个引理解释了为什么代价函数的 Monge 性足以传递到动态规划矩阵，并进一步给出每个候选前缀上的阶梯全单调性。SMAWK 适用于离线矩阵行最值搜索；但在式 \eqref{eq:basic-dp} 中，$A_{i,j}$ 又包含尚未全部求出的 $dp_j$，矩阵元素只能按下标从小到大的顺序逐步变得可用。这就是在线矩阵搜索与普通离线矩阵搜索的区别，也是 LARSCH 相比直接 SMAWK 更自然的地方\cite{larmore1991online,noshi2021larschlib}。

\subsection{与常见竞赛优化的关系}

常见竞赛模板之间的关系可用 \figref{fig:landscape} 概括。分治优化只需要最优决策单调，因此要求较弱，但通常是 $\BigO(n\log n)$ 或带额外层数；SMAWK 要求全单调性，能在线下矩阵中做到线性；LARSCH 面向在线 DP，理论上也可做到线性。简化 LARSCH 则牺牲一个对数因子，换取更低实现难度和更好的滑动代价兼容性。

\begin{figure}[H]
\centering
\begin{tikzpicture}[
  node distance=9mm and 8mm,
  box/.style={draw=topazsection, rounded corners, align=center, inner sep=5pt, fill=blue!3},
  arrow/.style={-Latex, thick, color=topazsection}
]
\node[box] (monge) {代价函数 $w(j,i)$\\满足 Monge 性};
\node[box, below=of monge] (tm) {隐式矩阵 $A_{i,j}$\\具有全单调性};
\node[box, below left=of tm] (offline) {离线行最值\\SMAWK};
\node[box, below=of tm] (online) {在线 DP\\LARSCH};
\node[box, below right=of tm] (simple) {竞赛实现\\简化 LARSCH};
\node[box, below=of online] (dp) {$dp_i=\min_{j<i}\{dp_j+w(j,i)\}$\\近线性或线性求解};
\draw[arrow] (monge) -- (tm);
\draw[arrow] (tm) -- (offline);
\draw[arrow] (tm) -- (online);
\draw[arrow] (tm) -- (simple);
\draw[arrow] (offline) |- (dp);
\draw[arrow] (online) -- (dp);
\draw[arrow] (simple) |- (dp);
\end{tikzpicture}
\caption{Monge 性、全单调矩阵搜索与在线 DP 优化的关系。}
\label{fig:landscape}
\end{figure}

\section{历史代表结果}
\label{sec:history}

\tabref{tab:history} 按时间列出 Monge DP 优化的代表结果。近期竞赛材料中的许多技巧并不是孤立出现的，而是把 20 世纪 70 年代到 90 年代的动态规划加速理论重新用于现代 OJ 题。

\begin{table}[htbp]
\centering
\caption{Monge DP 优化相关代表结果。}
\label{tab:history}
\small
\begin{tabular}{p{0.14\textwidth}p{0.26\textwidth}p{0.46\textwidth}}
\toprule
时间 & 代表工作 & 作用 \\
\midrule
1971 & Knuth 最优二叉搜索树\cite{knuth1971obst} & 利用最优根的单调性，将一类区间 DP 从三次降到二次。 \\
1980 & Yao 四边形不等式\cite{yao1980efficient} & 给出更一般的 quadrangle inequality 条件，解释何时可做 Knuth-Yao 优化。 \\
1987 & SMAWK\cite{aggarwal1987geometric} & 对全单调矩阵做隐式行最值搜索，是 Monge DP 离线优化的核心工具。 \\
1991 & Larmore-Schieber 在线 DP\cite{larmore1991online} & 将全单调矩阵搜索扩展到在线动态规划，并应用于 RNA 二级结构预测。 \\
1992 & 稀疏与凸性 DP\cite{eppstein1992sparse1,eppstein1992sparse2,galil1992dynamic} & 系统研究稀疏、凸性和凹性结构下的动态规划加速，展示 Monge 型结构与序列优化之间的联系。 \\
1996 & Monge 性综述\cite{burkard1996perspectives} & 从优化角度综述 Monge 性在排序、匹配、运输和矩阵搜索等问题中的作用。 \\
2009 & Knuth-Yao 与全单调性的统一\cite{bein2009knuthyao} & 证明 Knuth-Yao 加速可视为全单调性的结果，并给出在线版本解释。 \\
2021-2025 & 竞赛模板化与再传播\cite{noshi2021larschlib,noshi2023simplifiedlarsch,luogu2025larsch} & 将 LARSCH 和简化 LARSCH 写成可复用 C++ 模板，用于 WQS 二分、滑动窗口代价与莫队代价。 \\
\bottomrule
\end{tabular}
\end{table}

从算法发展脉络看，历史工作、理论定义、现代竞赛应用之间可以自然衔接。Knuth-Yao 优化解释“为什么决策单调有用”，SMAWK 解释“全单调矩阵可以怎样搜索”，LARSCH 则解释“在线 DP 为什么比离线矩阵搜索难”。除了竞赛中的模板化传播，Monge 矩阵和 totally monotone matrix searching 在计算几何、序列优化、最短路、RNA 二级结构预测等问题中也长期作为基础工具出现\cite{burkard1996perspectives,eppstein1992sparse1,eppstein1992sparse2,galil1992dynamic,larmore1991online}。这里重点放在在线 DP 与竞赛实现，因此不展开所有应用；这些正式文献说明，这条技术线并非只服务于 OJ 题，而是有更广的算法背景。


\section{一个核心算法：简化 LARSCH 分治框架}
\label{sec:algorithm}

\subsection{算法目标与假设}

完整 LARSCH 可以在线性时间求解在线 Monge DP，但实现较复杂。竞赛中更常见的是简化 LARSCH：多一个对数因子，代码明显更短，也更容易和滑动窗口或莫队维护的代价函数结合\cite{noshi2023simplifiedlarsch,luogu2025larsch}。

继续考虑式 \eqref{eq:basic-dp}。为了让证明条件精确，记
\[
  q_i^{(t)}=\argmin_{0\le j\le t,\,j<i}\{dp_j+w(j,i)\}
\]
为第 $i$ 行在候选前缀 $0,1,\ldots,t$ 中的最左最优列。算法需要的不是单纯的全局决策单调，而是每个候选前缀上的阶梯全单调性：对任意固定 $t$，只要 $i_1<i_2$ 且两行都存在合法候选，就有 $q_{i_1}^{(t)}\le q_{i_2}^{(t)}$。在常见 Monge 代价下，列平移引理和引理 \ref{lem:staircase-prefix} 给出了这个条件。

\subsection{伪代码}

以下骨架可直接翻译成 C++ 模板。状态为 $0,1,\ldots,n$，另设终点哨兵 $N=n+1$，因此数组长度必须为 $N+1=n+2$。函数 \lstinline{relax} 只在严格变优时更新，平局保留更早写入的决策，从而固定为“取最左最优”；实际实现还应按题目数据范围设置足够大的 \lstinline{INF}，并避免加法溢出。

\begin{minted}{cpp}
const long long INF = (1LL << 62);
int N = n + 1;                 // sentinel target
vector<long long> dp(N + 1, INF);
vector<int> pre(N + 1, 0);

void relax(int i, int j) {
    if (j >= i || dp[j] == INF) return;
    long long cand = dp[j] + cost(j, i);
    if (cand < dp[i]) {        // keep the leftmost optimum on ties
        dp[i] = cand;
        pre[i] = j;
    }
}

void solve(int l, int r) {
    if (r - l == 1) return;
    int mid = (l + r) >> 1;

    for (int j = pre[l]; j <= pre[r]; ++j) {
        relax(mid, j);
    }
    solve(l, mid);

    for (int j = l + 1; j <= mid; ++j) {
        relax(r, j);
    }
    solve(mid, r);
}

dp[0] = 0;
pre[0] = 0;
relax(N, 0);                   // only if cost(0,N) is defined
solve(0, N);
\end{minted}

哨兵 $N$ 的含义由具体问题决定：在最短路式 DP 中，它通常表示终点。若原模型不存在直接边 $(0,N)$，证明中可把它建模为代价 $+\infty$ 的形式候选，并令 $pre[N]=0$ 表示最左的无穷候选；实际代码应使用单独初始化，或在 \lstinline{relax} 中跳过这条边，而不是调用未定义的 \lstinline{cost(0,N)}。关键是保持同样的不变量语义，防止无穷值参与加法溢出，并保证所有有效更新都对应题目中的合法代价。

\subsection{正确性证明}

\begin{lemma}[递归不变量]
调用 \lstinline{solve(l, r)} 时满足三条性质：
\begin{enumerate}[(i)]
  \item $dp_0,dp_1,\ldots,dp_l$ 都已经是最终值；
  \item $dp_r$ 已经等价地取得候选前缀 $0,1,\ldots,l$ 下的最小值，且 $pre[r]=q_r^{(l)}$；
  \item $pre[l]=q_l^{(l)}$。当 $l=0$ 时没有合法前驱，约定 $pre[0]=0$。
\end{enumerate}
\end{lemma}

\begin{proof}
初始调用为 \lstinline{solve(0, N)}。此时只有 $dp_0=0$ 需要被视为最终值；若直接哨兵边不存在，则把候选 $0\to N$ 的代价视为 $+\infty$ 并保留 $pre[N]=0$；否则先用候选 $0$ 更新 $N$。两种情况下 $N$ 都已经取得候选前缀 $0$ 下的最小值，并有 $pre[N]=q_N^{(0)}$，不变量成立。

设某次调用 \lstinline{solve(l, r)} 满足不变量，令 $mid=(l+r)/2$。由于 $dp_0,\ldots,dp_l$ 都已最终，可以安全计算这些候选对 $mid$ 的贡献；又由 $pre[r]=q_r^{(l)}\le l$ 可知第一段循环不会访问尚未最终确定的 $dp_j$。在候选前缀 $0,1,\ldots,l$ 下，阶梯形全单调性给出
\[
  q_l^{(l)}\le q_{mid}^{(l)}\le q_r^{(l)}.
\]
其中 $q_l^{(l)}=pre[l]$，$q_r^{(l)}=pre[r]$。第一段循环按递增顺序枚举 $pre[l],pre[l]+1,\ldots,pre[r]$，因此一定会枚举到 $q_{mid}^{(l)}$；严格变优才更新的平局规则保留该范围内最左最优列，从而得到 $pre[mid]=q_{mid}^{(l)}$。于是左递归 \lstinline{solve(l, mid)} 的不变量成立。

左递归结束后，$dp_0,dp_1,\ldots,dp_{mid}$ 都已经成为最终值。第二段循环按递增顺序把新候选 $l+1,l+2,\ldots,mid$ 用于更新 $r$；旧候选前缀 $0\ldots l$ 已在调用开始前处理过，新候选下标都更大，所以相同代价下保留旧的更左决策仍符合最左规则。循环结束后，$r$ 正好取得候选前缀 $0\ldots mid$ 下的最小值，并且 $pre[r]=q_r^{(mid)}$。左递归结束也使 $mid$ 成为最终状态；由于 $j=mid$ 对第 $mid$ 行非法，$q_{mid}^{(mid)}=q_{mid}^{(mid-1)}=pre[mid]$。于是右递归 \lstinline{solve(mid, r)} 的不变量也成立。
\end{proof}

\begin{theorem}[简化 LARSCH 的正确性]
若隐式矩阵在合法阶梯区域上满足上述全单调条件，且 \lstinline{relax(i, j)} 正确计算候选 $dp_j+w(j,i)$，则算法结束后每个被求解状态的 $dp$ 值均为式 \eqref{eq:basic-dp} 的最优值。
\end{theorem}

\begin{proof}
考虑任意叶子调用 \lstinline{solve(k-1, k)}。由递归不变量，调用该叶子时 $dp_k$ 已经被所有 $0\le j\le k-1$ 的候选更新过，而这些正是式 \eqref{eq:basic-dp} 的全部合法候选。因此 $dp_k$ 已经达到最优。所有状态都会作为某个叶子的右端点出现，结论成立。
\end{proof}

\subsection{复杂度分析}

递归树每一层中，第二段循环的区间 $[l+1,mid]$ 对不同节点互不相交，总长度为 $\BigO(n)$。第一段循环的区间也可以线性求和：同一深度的递归节点按端点排成 $x_0<x_1<\cdots<x_s$，各节点枚举的列区间为 $[pre[x_a],pre[x_{a+1}]]$。这里的 $pre[x_a]$ 是递归不变量下已经得到的候选前缀最左最优列，不必等同于最终全局最优列。由阶梯全单调性，$pre[x_0],pre[x_1],\ldots,pre[x_s]$ 非降；并且这里的集合扩张始终是候选前缀扩张，新增候选列下标更大，平局又保留最左候选，所以最左最优列只会保持不变或右移。因此这些长度相加至多为 $pre[x_s]-pre[x_0]+s=\BigO(n)$。递归深度为 $\BigO(\log n)$，若 \lstinline{relax} 为 $\BigO(1)$，总复杂度为
\[
  T(n)=2T(n/2)+\BigO(n)=\BigO(n\log n).
\]

若代价 $w(j,i)$ 只能用滑动窗口维护，还需要把每层查询组织成两组单调移动的序列：第一段循环服务于中点查询，第二段循环服务于右端点查询。两组状态分开维护时，每层指针移动仍可控制在 $\BigO(n)$；若把两种方向混在一个窗口里，复杂度分析通常不再成立。这也是洛谷文章强调 CF868F 一类问题要拆分查询状态的原因\cite{luogu2025larsch}。

\section{算法竞赛中的应用}
\label{sec:applications}

\subsection{WQS 二分与 Aliens trick}

很多分段 DP 写作
\[
  dp_{i,k}=\min_{j<i}\{dp_{j,k-1}+w(j,i)\},
\]
其中 $k$ 是段数。WQS 二分，也常称 Aliens trick，会把段数限制用惩罚参数 $\lambda$ 吸收到目标函数中，变成
\[
  f_i(\lambda)=\min_{j<i}\{f_j(\lambda)+w(j,i)+\lambda\}.
\]
这样每次固定 $\lambda$ 都只需求解一维在线 Monge DP。实现中通常同时维护二元组“代价、段数”；一种常用约定是在惩罚后总代价相同的时候取更多段；随着 $\lambda$ 增大，记录的段数不增。二分时选择让最优段数跨过目标值的边界参数，再用惩罚后代价减去 $\lambda$ 乘目标段数恢复原目标；若边界处记录的段数并不恰好等于目标段数，恢复值可理解为离散凸包上该斜率支持线在目标段数处的截距。AtCoder ABC355 G 官方题解明确列出：分治 + monotone minima 为 $\BigO(N(\log N)^2)$，分治 + SMAWK 为 $\BigO(N\log N)$，LARSCH 为 $\BigO(N)$，同时也给出简化 LARSCH 在实现常数上的优势\cite{atcoder2024abc355g}。


\subsection{完整案例：ABC355G Baseball}

同样的判断流程可以在 AtCoder ABC355 G Baseball 中完整落地。题目给定权重序列 $P_1,\ldots,P_N$，先选择 $K$ 个不同位置，再按 $P_y$ 成比例随机产生位置 $y$，目标是最小化 $y$ 到最近被选位置的期望距离乘以 $\sum_y P_y$；约束为 $N\le 5\times 10^4$，朴素二次或多维 DP 都不可接受\cite{atcoder2024abc355g-task}。

官方题解先把问题改写成有向无环图最短路\cite{atcoder2024abc355g}。图的顶点为 $0,1,\ldots,N,N+1$，其中 $0$ 和 $N+1$ 是左右边界哨兵。边权定义为
\begin{align}
  c(0,j)&=\sum_{y=1}^{j-1} P_y(j-y), &&1\le j\le N, \\
  c(i,j)&=\sum_{y=i}^{j-1} P_y\min(y-i,j-y), &&1\le i<j\le N, \\
  c(i,N+1)&=\sum_{y=i}^{N} P_y(y-i), &&1\le i\le N .
  \label{eq:baseball-cost}
\end{align}
原图没有直接边 $(0,N+1)$；若为了统一模板形式引入这个位置，也应令 $c(0,N+1)=+\infty$，保证它不会参与最优路径。若选择位置为 $x_1<\cdots<x_K$，则路径 $0\to x_1\to\cdots\to x_K\to N+1$ 的权重正好等于目标值。因此问题等价于求从 $0$ 到 $N+1$、恰好使用 $K+1$ 条边的最短路。令 $D_{t,j}$ 表示使用 $t$ 条边到达 $j$ 的最小代价，就得到
\[
  D_{t,j}=\min_{i<j,\,(i,j)\in E}\{D_{t-1,i}+c(i,j)\},
\]
其中 $E$ 是上面三类边的集合。朴素计算为 $\BigO(N^2K)$。

这些边权可以用 $P_y$ 与 $P_y y$ 的前缀和在 $\BigO(1)$ 时间求出。例如令 $S(l,r)=\sum_{y=l}^rP_y$、$T(l,r)=\sum_{y=l}^rP_yy$，则
\[
  c(0,j)=jS(1,j-1)-T(1,j-1),\qquad
  c(i,N+1)=T(i,N)-iS(i,N),
\]
而对 $1\le i<j\le N$，取 $m=\lfloor(i+j)/2\rfloor$，有
\[
  c(i,j)=T(i,m)-iS(i,m)+jS(m+1,j-1)-T(m+1,j-1),
\]
空区间贡献按 $0$ 处理。

WQS 二分把边数限制吸收到惩罚参数 $\lambda$ 中：每条边额外加上 $\lambda$，反复求解
\begin{equation}
  F_j(\lambda)=\min_{i<j,\,(i,j)\in E}\{F_i(\lambda)+c(i,j)+\lambda\},
  \label{eq:baseball-aliens}
\end{equation}
并同时记录最短路使用的边数。实现上同样维护“代价、边数”二元组；可约定惩罚后总代价相同则取更多边，使记录的边数随 $\lambda$ 增大而不增。题目数据为整数时，可以在整数 $\lambda$ 上二分搜索使最优边数仍不少于 $K+1$ 的边界参数，或等价地搜索从“不少于”变为“少于”的临界点；在该边界处用 $F_{N+1}(\lambda)-\lambda(K+1)$ 恢复恰好 $K+1$ 条边的原目标值。即使边界最短路本身的记录边数不是 $K+1$，离散凸性也保证这个支持线截距给出目标边数处的最优值。官方题解给出的搜索范围为 $\BigO(NS)$ 量级，其中 $S=\sum_i P_i$\cite{atcoder2024abc355g}。

剩下的关键是证明边权满足 Monge 性。由于 $P_y$ 是非负权重，可以先对每个位置 $y$ 逐点比较贡献，再对 $y$ 求和。对四条相关边均存在的端点 $0\le a<b<c<d\le N+1$，需要
\begin{equation}
  c(a,c)+c(b,d)\le c(a,d)+c(b,c).
  \label{eq:baseball-monge}
\end{equation}
若形式上出现 $(0,N+1)$，则 $c(0,N+1)=+\infty$ 位于右侧，式 \eqref{eq:baseball-monge} 平凡成立；下面只需检查四条边均有限的情况。

先考虑四个端点均在 $1,\ldots,N$ 内的情形。固定一个位置 $y$，只比较它对式 \eqref{eq:baseball-monge} 两侧的贡献。若 $a\le y<b$，右侧比左侧多
\[
  \min(y-a,d-y)-\min(y-a,c-y)\ge 0.
\]
若 $b\le y<c$，令 $A=y-a$、$B=y-b$、$C=c-y$、$D=d-y$。有 $A\ge B$ 且 $D\ge C$。函数 $h(t)=\min(t,D)-\min(t,C)$ 关于 $t$ 单调非降，因此
\[
  \min(A,D)+\min(B,C)\ge \min(A,C)+\min(B,D).
\]
若 $c\le y<d$，右侧比左侧多
\[
  \min(y-a,d-y)-\min(y-b,d-y)\ge 0.
\]
其他位置贡献为零。

边界 $a=0$、$d\le N$ 时也可逐点比较右侧减左侧。若 $1\le y<b$，差为 $d-c\ge0$；若 $b\le y<c$，差为
\[
  d-c+\min(y-b,c-y)-\min(y-b,d-y)\ge0,
\]
因为把第二个参数从 $c-y$ 增加到 $d-y$ 时，$\min(y-b,\cdot)$ 的增量不超过 $d-c$；若 $c\le y<d$，差为
\[
  d-y-\min(y-b,d-y)\ge0.
\]

边界 $a\ge1$、$d=N+1$ 时同理。若 $a\le y<b$，右侧减左侧为
\[
  y-a-\min(y-a,c-y)\ge0;
\]
若 $b\le y<c$，差为
\[
  b-a+\min(y-b,c-y)-\min(y-a,c-y)\ge0,
\]
因为 $y-a=(y-b)+(b-a)$，而 $\min(\cdot,c-y)$ 的增量不超过 $b-a$；若 $c\le y\le N$，差为 $b-a\ge0$。乘以非负权重 $P_y$ 并对所有 $y$ 求和，即得式 \eqref{eq:baseball-monge}。由此，固定 $\lambda$ 后的式 \eqref{eq:baseball-aliens} 正好落入在线 Monge DP 框架。

\tabref{tab:abc355g-pipeline} 总结了这个案例的解题链条。关键步骤不是机械套用 LARSCH 模板，而是先将题意转成图模型，再证明边权结构，最后根据在线依赖选择合适的矩阵搜索算法。

\begin{table}[htbp]
\centering
\caption{ABC355G Baseball 的 Monge/LARSCH 建模链条。}
\label{tab:abc355g-pipeline}
\begin{tabular}{p{0.22\textwidth}p{0.64\textwidth}}
\toprule
步骤 & 内容 \\
\midrule
题意压缩 & 选择 $K$ 个位置，最小化带权距离到最近选择点的总和。 \\
图模型 & 相邻选择点 $(i,j)$ 形成一条边，边权为区间内到较近端点的加权距离；左右哨兵使用单侧距离，且没有有效边 $(0,N+1)$。 \\
段数约束 & 恰好 $K+1$ 条边的最短路，用 WQS 二分转为每条边加惩罚 $\lambda$。 \\
Monge 证明 & 四个端点 $a<b<c<d$ 下逐点比较贡献，得到 $c(a,c)+c(b,d)\le c(a,d)+c(b,c)$。 \\
在线依赖 & $F_j$ 依赖更小下标的 $F_i$，不能把整张隐式矩阵直接交给离线 SMAWK。 \\
实现选择 & 前缀和使边权 $\BigO(1)$ 查询；简化 LARSCH 每次评估 $\lambda$ 为 $\BigO(N\log N)$，总复杂度为 $\BigO(N\log N\log(NS))$。 \\
\bottomrule
\end{tabular}
\end{table}

\subsection{滑动窗口代价与莫队代价}

有些转移代价 $w(l,r)$ 不容易用闭式公式 $\BigO(1)$ 计算，但可以在移动左右端点时维护。例如区间内相等数对数量、合法括号子串数量等。这时，普通分治优化虽然知道候选区间，却不一定能安排代价查询顺序；SMAWK 要求随机访问隐式矩阵元素，也会破坏滑动窗口维护。简化 LARSCH 的优势在于递归每层可以把查询组织成两套单调移动的窗口。洛谷的 2025 年文章特别强调，若使用莫队维护代价，第一步和第三步应分开维护两个窗口，否则指针会来回跳动导致复杂度退化\cite{luogu2025larsch}。

\subsection{min-plus/max-plus 卷积与背包优化}

另一个常见入口是卷积。给定两个数组 $A,B$，max-plus 卷积定义为
\[
  C_i=\max_{j+k=i}(A_j+B_k).
\]
一般情形很难突破二次复杂度，但如果至少一边满足与 min/max 方向相容的凸性或凹性，问题可以转化为 Monge 矩阵行最值搜索，使用 SMAWK 或斜率合并得到线性或近线性算法。Codeforces 上关于背包、子集和与 max-plus convolution 的教程明确将“凸数组上的 max-plus convolution 可由 SMAWK 线性求解”列为核心工具\cite{codeforces2022knapsack}。2026 年 Codeforces 上基于 THUPC 的题解也展示了把 DP 更新理解为 min-plus convolution，并利用凸函数斜率集合压缩状态的思路\cite{codeforces2026thupc}。

\subsection{与 C++ 模板化的关系}

从 C++ 竞赛实践看，这类算法的真正门槛往往不是定理本身，而是三个工程问题：
\begin{enumerate}[(1)]
  \item 如何证明题目中的 $w$ 确实满足 Monge 性；
  \item 如何在不破坏复杂度的顺序下查询 $w$；
  \item 如何在平局时稳定选择最左或最右决策，保证单调性。
\end{enumerate}
因此，竞赛模板不应只保存代码，还应保存“证明检查表”。对应检查表见 \tabref{tab:checklist}。

\begin{table}[H]
\centering
\small
\renewcommand{\arraystretch}{0.92}
\caption{在竞赛中判断是否可用 Monge/LARSCH 的检查表。}
\label{tab:checklist}
\begin{tabular}{p{0.22\textwidth}p{0.64\textwidth}}
\toprule
检查项 & 需要确认的内容 \\
\midrule
转移形式 & 是否能写成 $dp_i=\min_{j<i}\{dp_j+w(j,i)\}$，或经参数化后化为同形。 \\
Monge 证明 & 是否能证明四点不等式，且方向、下标顺序与题目代价一致。 \\
在线依赖 & 若 $A_{i,j}$ 含尚未计算的 $dp$，就不能直接离线 SMAWK。 \\
代价访问 & $w$ 是 $\BigO(1)$、前缀和、滑动窗口还是莫队；查询顺序是否可维护。 \\
平局规则 & 最优决策固定取最左或最右，并与单调性证明一致。 \\
实现选择 & 离线矩阵用 SMAWK；在线 DP 用简化 LARSCH；特化函数族可用凸包或 Li Chao tree。 \\
\bottomrule
\end{tabular}
\end{table}

\section{State-of-the-art、局限与实践判断}
\label{sec:sota}


\subsection{理论 SOTA 与实践 SOTA}

这里的 state-of-the-art 需要分成两个语境。\textbf{理论 SOTA} 关注渐进复杂度和矩阵查询次数：对离线 totally monotone 矩阵，SMAWK 已达到线性数量级查询\cite{aggarwal1987geometric}；对在线 Monge DP，完整 LARSCH 也达到 $\BigO(n)$\cite{larmore1991online,noshi2021larschlib}。从这个角度看，线性查询或线性时间是最清楚的理论基准。

\textbf{实践 SOTA} 则不同。算法竞赛中的“最优”还要考虑代码长度、常数、调试难度、平局规则以及代价函数能否按指定顺序访问。完整 LARSCH 虽然渐进最优，但手写难度高；简化 LARSCH 牺牲一个对数因子，得到更短的代码和更稳定的查询顺序，尤其适合滑动窗口或莫队代价。AtCoder 官方题解也提到，虽然 LARSCH 有 $\BigO(N)$ 解法，但简化版 $\BigO(N\log N)$ 在作者实现中常数更低，实际运行时间可能更好\cite{atcoder2024abc355g}。

\tabref{tab:sota-compare} 把几种常见方案放在同一张表中。理论最优关注矩阵查询次数或单次在线 DP 的渐进复杂度；竞赛实现则要按代价访问方式在简化 LARSCH、SMAWK、monotone minima 与特化凸包之间取舍。

\begin{table}[htbp]
\centering
\caption{Monge DP 相关方法的理论与实践对照。}
\label{tab:sota-compare}
\begin{tabular}{p{0.18\textwidth}p{0.24\textwidth}p{0.20\textwidth}p{0.24\textwidth}}
\toprule
方法 & 适用条件 & 复杂度 & 实现风险 \\
\midrule
分治优化 & 全局决策单调，按层离线求值 & 约 $\BigO(n\log n)$ & 条件较弱，但无法处理强在线依赖。 \\
Monotone minima & 行最优列单调，可随机访问元素 & 约 $\BigO(n\log n)$ & 代码简单，常数小；随机访问可能破坏滑动代价。 \\
SMAWK & 全单调矩阵，可随机访问元素 & $\BigO(n)$ 查询 & 渐进最优，代码和调试成本较高。 \\
完整 LARSCH & 在线 DP 的阶梯全单调矩阵 & $\BigO(n)$ & 理论最优，实现复杂，竞赛中较少手写。 \\
简化 LARSCH & 在线 DP，允许一个对数因子 & $\BigO(n\log n)$ & 代码较短，适合与滑动窗口或莫队代价结合。 \\
凸包/Li Chao & 函数交点单调或线性函数族 & 通常 $\BigO(n\log n)$ & 依赖题目特化结构，不是通用 Monge 模板。 \\
\bottomrule
\end{tabular}
\end{table}

\subsection{主要局限}

这一技术线也有明显局限。

\textbf{第一，Monge 性证明高度依赖题目结构。} 很多题并不会直接告诉选手 $w$ 满足四边形不等式，需要通过几何、排序、交换论证或凸性推导得到。一旦证明方向错误，仅掌握模板也无法保证正确性。

\textbf{第二，一般 min-plus convolution 不在这条模板线内。} Monge/SMAWK 的加速依赖特殊结构。若两个数组都没有凸性、凹性或类似单调结构，就不能直接套用本综述讨论的线性或近线性矩阵搜索技巧。

\textbf{第三，工程实现容易被访问顺序破坏。} 如果 $w$ 只能滑动维护，随机访问会让本来正确的渐进分析失效。洛谷文章中对 CF868F 的提醒就是典型例子：两个方向的查询必须分开维护，不能共用同一个莫队状态\cite{luogu2025larsch}。


\subsection{本文的综合诠释：在线 Monge DP 的四层判定框架}

本文的一个综合性观点是：Monge DP 优化不应被理解为若干孤立模板，而应被理解为“结构证明—矩阵搜索—在线依赖—代价访问”的四层判定流程。这类方法的价值不只在于某个模板本身，而在于提供了一套可复用的判断流程：
\begin{enumerate}[(1)]
  \item 先把“决策单调性 DP 优化”改写成“隐式矩阵行最值搜索”；
  \item 再区分 SMAWK 与 LARSCH 的核心差异：前者要求离线随机访问，后者处理在线依赖；
  \item 然后根据代价访问方式选择完整 LARSCH、简化 LARSCH、monotone minima 或特化凸包；
  \item 最后用 \tabref{tab:checklist} 与 \tabref{tab:abc355g-pipeline} 这样的检查表，把证明条件和实现风险逐项核对。
\end{enumerate}

在 ACM/OJ 的语境下，这种解释比直接背诵 SMAWK 的 reduce 步骤更贴近解题过程。熟悉的 $dp_i=\min_j$ 转移只是入口；更关键的是候选区间为什么可以缩小、在线依赖为什么阻止直接离线搜索、以及代价函数能否按所需顺序被高效访问。


\section{结论}

Monge 性、SMAWK 与 LARSCH 并不是 2025-2026 年才出现的新算法；近期竞赛题解和模板文章中的新价值，主要体现在这些经典工具被重新组合到在线 DP、WQS 二分和滑动代价维护之中。单纯掌握“分治优化 DP”不足以覆盖这类题目；更关键的是 Monge 不等式能否成立、在线依赖如何处理、代价函数能否按所需顺序访问。

ABC355G 案例说明，Monge 优化不是孤立模板，而是从题意建模、结构证明到实现取舍的一套分析流程。竞赛中看似经验化的优化，往往可以回到 Monge 不等式和隐式矩阵搜索来解释；把在线依赖、代价访问顺序和平局规则合在同一张检查表中，也使 SMAWK、LARSCH 与简化实现之间的取舍具有清晰的数学依据。

\clearpage
\printbibliography[heading=bibintoc, title=\ebibname]

\end{document}
